\chapter{VAR and Granger Causality}
\label{cap:var}
% ── index ───────────────────────────────────────────────────────
\index{VAR|textbf}
\index{vector autoregressive|see{VAR}}
\index{var@\texttt{var()}|textbf}
\index{Granger causality|textbf}
\index{granger@\texttt{granger()}|textbf}
\index{impulse response function|textbf}
\index{irf@\texttt{irf()}|textbf}
\index{variance decomposition}

% ─────────────────────────────────────────────────────────────────────────────
\section{The VAR model}

The Vector Autoregressive model (\textbf{VAR($p$)}) generalizes AR to
multivariate systems. With $K$ variables and $p$ lags:

\[
  \mathbf{y}_t = \mathbf{c} + \mathbf{A}_1 \mathbf{y}_{t-1}
                 + \cdots + \mathbf{A}_p \mathbf{y}_{t-p}
                 + \boldsymbol{\varepsilon}_t
\]

where $\mathbf{y}_t = (y_{1t}, \ldots, y_{Kt})'$ is the vector of variables,
$\mathbf{A}_j$ are $K \times K$ coefficient matrices and
$\boldsymbol{\varepsilon}_t \sim (0, \Sigma_u)$ is the vector of shocks that are
uncorrelated over time but potentially correlated with each other.

The VAR is estimated equation by equation via OLS --- each equation is an AR($p$)
in the $K$ lagged variables. The stationarity condition requires that all
eigenvalues of $\mathbf{A}_1 + \cdots + \mathbf{A}_p$ lie inside the
unit circle.

% ─────────────────────────────────────────────────────────────────────────────
\section{Granger Causality}

$x$ \textbf{Granger-causes} $y$ if lags of $x$ have predictive power
for $y$ beyond what lags of $y$ alone provide. Formally, the hypothesis tested is:
\[
  H_0: A_{12}(1) = A_{12}(2) = \cdots = A_{12}(p) = 0
\]
in the equation for $y_1$, where $A_{12}(j)$ is the coefficient of $y_{2,t-j}$.

The test is a standard F-test in the auxiliary equation:
\[
  F = \frac{(RSS_r - RSS_u)/p}{RSS_u / (T - 2p - 1)} \sim F(p,\; T-2p-1)
\]

\begin{nota}
  Granger causality is \emph{predictability}, not structural causality.
  $x$ that precedes $y$ and correlates with it passes the test,
  even without a direct causal relationship.
\end{nota}

% ─────────────────────────────────────────────────────────────────────────────
\section{Impulse Response Function}

The \textbf{IRF} (\emph{Impulse Response Function}) traces the effect of a
unit shock in $y_k$ on all variables over time.

The MA($\infty$) representation of the VAR is $\mathbf{y}_t = \sum_{h=0}^\infty
\mathbf{\Phi}_h \boldsymbol{\varepsilon}_{t-h}$, with $\mathbf{\Phi}_0 = I$
and $\mathbf{\Phi}_h = \mathbf{A}_1 \mathbf{\Phi}_{h-1} + \cdots +
\mathbf{A}_p \mathbf{\Phi}_{h-p}$. The orthogonalized IRF uses the Cholesky
decomposition of $\Sigma_u$ to isolate uncorrelated shocks --- the
\textbf{Cholesky ordering} (variable ordering) matters.

% ─────────────────────────────────────────────────────────────────────────────
\section{DGP Verification}

The DGP uses a bivariate VAR(1) with an asymmetric causality structure:
$y_2$ causes $y_1$ (coefficient $0{,}2$), but $y_1$ causes $y_2$
weakly (coefficient $0{,}1$):

\[
  \begin{pmatrix} y_{1t} \\ y_{2t} \end{pmatrix}
  = \begin{pmatrix} 0{,}5 & 0{,}2 \\ 0{,}1 & 0{,}6 \end{pmatrix}
    \begin{pmatrix} y_{1,t-1} \\ y_{2,t-1} \end{pmatrix}
  + \begin{pmatrix} \varepsilon_{1t} \\ \varepsilon_{2t} \end{pmatrix}
\]

\begin{lstlisting}[caption={VAR(1) DGP --- Granger and IRF}]
seed(42)
input df
id
1
end
for i in 1..=9 { df = append(df, df) }
df |> filter(_n <= 300)
generate df e1 = rnormal()
generate df e2 = rnormal()
generate df y1 = 0.0
generate df y2 = 0.0
generate df t  = _n
let n = nrow(df)
let i = 2
while i <= n {
    let l1 = mean(df, y1, if = t == i - 1)
    let l2 = mean(df, y2, if = t == i - 1)
    let a  = mean(df, e1, if = t == i)
    let b  = mean(df, e2, if = t == i)
    replace df y1 = 0.5*l1 + 0.2*l2 + a if t == i
    replace df y2 = 0.1*l1 + 0.6*l2 + b if t == i
    i = i + 1
}
let m = var(df, y1, y2, lags=1)
display m
granger(df, y1, y2, lags=1)
granger(df, y2, y1, lags=1)
irf(m, steps=6)
\end{lstlisting}

\subsection*{Model fit}

\begin{lstlisting}[language={}, basicstyle=\ttfamily\small,
  frame=none, numbers=none, backgroundcolor=\color{codbg},
  xleftmargin=0pt, xrightmargin=0pt, breaklines=false]
VAR(1)  N=299  AIC=-5.0269  BIC=-4.9526

Sigma_u:
  [ 0.0819  0.0046 ]
  [ 0.0046  0.0772 ]
\end{lstlisting}

The residual covariance matrix $\hat\Sigma_u$ is nearly diagonal --- the
shocks $\varepsilon_1$ and $\varepsilon_2$ are practically orthogonal,
as specified in the DGP.

\subsection*{Granger Causality}

\begin{lstlisting}[language={}, basicstyle=\ttfamily\small,
  frame=none, numbers=none, backgroundcolor=\color{codbg},
  xleftmargin=0pt, xrightmargin=0pt, breaklines=false]
================== Granger Causality Test ==================
  H₀: y2 does not Granger-cause y1   (lags=1)
  F(1, 296) = 9.4437   p = 0.0023
  Conclusion: REJECTS H₀ — y2 Granger-causes y1

================== Granger Causality Test ==================
  H₀: y1 does not Granger-cause y2   (lags=1)
  F(1, 296) = 4.5768   p = 0.0332
  Conclusion: REJECTS H₀ — y1 Granger-causes y2
\end{lstlisting}

Both directions reject $H_0$: $y_2 \to y_1$ with $p = 0{,}002$ and
$y_1 \to y_2$ with $p = 0{,}033$. The DGP imposes $\phi_{12} = 0{,}2$ and
$\phi_{21} = 0{,}1$ --- bidirectional causality is detected, although
$y_2 \to y_1$ is stronger.

\subsection*{Impulse response}

\begin{lstlisting}[language={}, basicstyle=\ttfamily\small,
  frame=none, numbers=none, backgroundcolor=\color{codbg},
  xleftmargin=0pt, xrightmargin=0pt, breaklines=false]
IRF — VAR(1) — 6 steps

  Impulse: y1
  ──────────────────────────────────────
       h          y1          y2
  ──────────────────────────────────────
       1      0.2862      0.0161
       2      0.1430      0.0382
       3      0.0758      0.0364
       4      0.0426      0.0286
       5      0.0251      0.0207
       6      0.0154      0.0145
  ──────────────────────────────────────

  Impulse: y2
  ──────────────────────────────────────
       h          y1          y2
  ──────────────────────────────────────
       1      0.0000      0.2774
       2      0.0405      0.1595
       3      0.0432      0.0958
       4      0.0352      0.0595
       5      0.0260      0.0378
       6      0.0183      0.0243
  ──────────────────────────────────────
\end{lstlisting}

A shock to $y_1$ has a rapid effect on $y_1$ (decays from $0{,}29$ to
$0{,}015$ over 6 periods) and a small effect on $y_2$. A shock to $y_2$
transmits progressively to $y_1$ --- reaches $0{,}043$ at period 3,
then decays --- capturing the $\phi_{12} = 0{,}2$ channel of the DGP.

% ─────────────────────────────────────────────────────────────────────────────
\section{Lag selection}

\begin{lstlisting}[caption={Comparison by number of lags}]
let m1 = var(df, y1, y2, lags=1)
let m2 = var(df, y1, y2, lags=2)
let m3 = var(df, y1, y2, lags=3)
esttab(m1, m2, m3)
\end{lstlisting}

The lag with the lowest BIC (parsimony) or AIC (fit) is selected.
With the VAR(1) DGP, $p=1$ is expected to minimize both.

% ─────────────────────────────────────────────────────────────────────────────
\section{Syntax}

\begin{lstlisting}
var(df, y1, y2, lags=1)
var(df, y1, y2, y3, lags=2)
granger(df, y, x, lags=1)
irf(m, steps=10)
\end{lstlisting}

\begin{lstlisting}[caption={Typical workflow --- GDP and inflation}]
load "macro.csv" as df

adf(df, gdp)
adf(df, cpi)

let m = var(df, gdp, cpi, lags=2)
granger(df, gdp, cpi, lags=2)
granger(df, cpi, gdp, lags=2)
irf(m, steps=12)
\end{lstlisting}

\begin{nota}
  VAR requires \emph{stationary} series. If the variables are I(1) and
  cointegrated, the appropriate model is the VECM (ch.~\ref{cap:coint}). If
  they are I(1) but not cointegrated, difference before estimating.
\end{nota}

\section{Empirical validation}

The VAR(1) DGP in this chapter corresponds to the \texttt{var\_book}
validation case in \texttt{validation/cases/}. The case compares
\hayashi{}'s equation-by-equation OLS estimation with R and Python
references. Run \lstinline{hay validate} to see the current status.
